\section{Thesis}

Thesis: \emph{Data-based Numerical Simulation of Buoyant Jet Flow}. Supervisor: Prof.\ Stefan
Pirker. Every number here is taken from the thesis of 10 October 2026. Equation and figure numbers
refer to that version; the derivations follow Appendix~D of the thesis.

\subsection{The talk in 15 minutes}

Build the talk on three claims, not on the chapters. One slide is about one minute.

\begin{center}
\begin{tabularx}{\linewidth}{@{}l l X@{}}
\toprule
Min & Slide & What you say \\
\midrule
0--1 & Title & Recurrence CFD records a flow once and replays it many times. I test it on a scalar
  that pushes back on the flow: temperature. \\
1--3 & Problem & rCFD was built for passive scalars. Temperature drives the flow by buoyancy. What
  does the replay lose, and how much can I repair without solving the flow again? \\
3--4 & Case & Closed tank, $75\times75\times15$~mm, 1~mm cells. Warm jet, 333~K at 0.02~m/s, for
  315~s. $Re \approx 99$, laminar. Two lids: adiabatic and cooled. \\
4--6 & Method & The preparation writes 600 frames. One replay step: recorded advection, face
  exchange, correction, conservation constraint. About 260 times cheaper. \\
6--8 & Claim 1 & No single mixing constant is best. The outlet sits just below the front, so a
  smeared front drains heat. \\
8--11 & Claim 2 & A bounded anti-diffusive correction, weighted to the vertical. Front 15.7~mm
  $\to$ 5.3 to 4.5~mm, reference 5.0~mm. Error of the mean temperature more than halved. \\
11--13 & Claim 3 & Under the cooled lid, on the fluid of the reference, the correction at
  $s \approx 0.55$ ends 0.06~K off on 9.62~K of cooling. \\
13--14 & What remains & While the tank fills, the layers separate too slowly. A flux read from the
  recording reduces this error from 2.25 to 0.45~K, but the tank then ends 0.86~K too cold. \\
14--15 & Take-home & A database belongs to a flow regime, not to a case. A numerical error at a
  front becomes an energy error near an outlet. \\
\bottomrule
\end{tabularx}
\end{center}

\begin{tipbox}
\textsf{The thesis in one sentence.} ``I tested recurrence CFD on temperature, a scalar that drives
its own flow. A bounded anti-diffusive correction removes the smearing of the front. The one defect
that is left is missing transport, and a flux read from the recording can supply it.''
\end{tipbox}

\subsection{Questions on the thesis}

\begin{qabox}{What problem did you solve?}
\origin{asked}{thesis part, committee Küng/Klambauer: ``What problem was solved?''}
\core{I tested recurrence CFD on temperature, a scalar that drives its own flow.}
\explain rCFD replays a recorded flow and moves a scalar with it. It was built for passive scalars,
which do not change the flow. Temperature does change the flow. So I measured what the replay loses
and how much I can repair without solving the flow again.
\nums{Reference 569~min. One replay 2.16~min on average. About 260 times faster.}
\summary{Passive scalar $\to$ active scalar. What is lost, what can be repaired.}
\src{Thesis, Abstract and Chapter 1.}
\end{qabox}

\begin{qabox}{Why should anybody care?}
\origin{asked}{thesis part, Mara (``What benefit does your thesis bring to the world?'') and Küng (``why should I care?'')}
\core{Thermal processes run for hours, and full CFD of that time is too slow for daily use.}
\explain rCFD is 260 times faster in my case. But a fast model is only useful if we know its error.
My thesis measures this error for a case where temperature matters, and it shows which part can be
repaired.
\summary{Fast is useful only with a known error.}
\src{Thesis, Chapter 1.}
\end{qabox}

\begin{qabox}{Explain recurrence CFD in simple words.}
\origin{predicted}{your thesis: everything you used can be asked}
\core{A normal CFD run is recorded once, and the replay reuses the recorded motion again and again.}
\explain The preparation solves the full equations and stores, frame by frame, which share of each
cell moves to which cell. The replay applies these stored maps to the temperature. Velocity and
pressure are never solved again, so it is cheap. One replay step has four stages.
\derive
\begin{align*}
\mathbf{T}^{*} &= \mathbf{M}^{(f_n)}\,\mathbf{T}^{n} \why{recorded map: move the heat} \\
\mathbf{T}^{**} &= \mathcal{S}(\mathbf{T}^{*}) \why{face exchange: add mixing} \\
\mathbf{T}^{***} &= \mathcal{C}(\mathbf{T}^{**}) \why{correction: my work} \\
\mathbf{T}^{n+1} &= \mathcal{B}(\mathbf{T}^{***}) \why{conservation constraint: fix the total}
\end{align*}
Together this is Eq.~(2.10), $\mathbf{T}^{n+1} = \mathcal{B}(\mathcal{C}(\mathcal{S}(\mathbf{M}\mathbf{T}^n)))$.
Standard rCFD has $\mathcal{C} = \text{identity}$.
\summary{Record once, replay many times. Map, swap, correct, balance.}
\src{Thesis, Sections 2.2--2.7, Eq.~(2.10).}
\end{qabox}

\begin{qabox}{Why is this AI, and not only CFD?}
\origin{asked}{thesis part, Küng committee (``Why and how does ML solve this problem?'')}
\core{It is a data-based simulation: an expensive solver produces data, and a cheap model reuses it.}
\explain The database plays the role of training data. Choosing the next frame is a
nearest-neighbour search in a database of states, the method of analogues. My main finding is about
generalization: a database belongs to a flow regime. A boundary condition that changes the regime
makes the database out of distribution.
\summary{Database = training data. New regime = out of distribution.}
\src{Thesis, Section 2.2 and Section 6.4.}
\end{qabox}

\begin{qabox}{What is the recurrence matrix, and is your flow recurrent?}
\origin{predicted}{your thesis: everything you used can be asked}
\core{The matrix compares every pair of recorded states, and my flow is not recurrent.}
\explain A small distance means two states are similar, so the replay could jump between them. In
my tank the matrix has two blocks, filling and settling, and no state comes back. So the replay
plays the frames in order. It is still rCFD, but the saving comes only from not solving velocity
and pressure.
\nums{Median prominence 0.028. The method asks for 0.05 to 0.1.}
\derive
\begin{align*}
R_{fg} &= \bigl\lVert \Phi^{(f)} - \Phi^{(g)} \bigr\rVert \why{distance between frames $f$ and $g$, Eq.~(2.5)} \\
d(f,g) &= R_{fg}/\max R \why{normalise to $[0,1]$, Eq.~(2.6)} \\
P_{\mathrm{rel}} &= \bigl\langle d_{\max} - d_{\min} \bigr\rangle / \bar{d} \why{depth of the dips over the plateau, Eq.~(2.7)}
\end{align*}
\summary{Two blocks, no return, 0.028 $\to$ frames in order.}
\src{Thesis, Section 2.4, Figure 2.2.}
\end{qabox}

\begin{qabox}{What is the face exchange, and what is $\fsd$?}
\origin{predicted}{your thesis: everything you used can be asked}
\core{After each advection the library moves a share of heat across every face, from warm to cold.}
\explain A pure displacement does not mix, so this step adds mixing. $\fsd$ is its strength per
step. It has no unit, so it is not a diffusion coefficient, but it can be translated into one. The
exchange can only reduce differences. It can never rebuild a smeared front.
\nums{$\fsd = 0.5$ gives $2.5\times10^{-7}$~m$^2$/s, ten times the molecular value.}
\derive
\begin{align*}
\Delta_{\mathrm{swap}} &= \fsd\,V\,\frac{T_{c_1}-T_{c_0}}{2} \why{exchange of Eq.~(2.11)} \\
Q_{\mathrm{swap}} &= \rho c_p\,\fsd\,c_s^{3}\,\frac{T_{c_1}-T_{c_0}}{2} \why{heat of that amount, $V = c_s^3$} \\
Q_{\mathrm{Fick}} &= \rho c_p\,D\,\frac{T_{c_1}-T_{c_0}}{c_s}\,c_s^{2}\,\Delta t \why{Fick's law over one face and one step} \\
D_{\mathrm{eq}} &= \frac{\fsd\,c_s^{2}}{2\,\Delta t} \why{set $Q_{\mathrm{Fick}} = Q_{\mathrm{swap}}$, Eq.~(2.13)}
\end{align*}
\summary{Swap = one-way mixing. Converts to $D_{\mathrm{eq}} = \fsd c_s^2/2\Delta t$.}
\src{Thesis, Section 2.6, Eqs.~(2.11)--(2.13); Appendix D.1.}
\end{qabox}

\begin{qabox}{Why is there no single best mixing constant?}
\origin{predicted}{your thesis: everything you used can be asked}
\core{The two measures want different amounts of mixing, because the outlet sits just below the front.}
\explain The mean temperature is best with little mixing. The height of the heat is best near
$\fsd = 0.5$. The more the replay smears the front, the more warm fluid reaches the outlet and
leaves the tank.
\nums{Mean-temperature error 1.07~K at little mixing. At $\fsd = 3$ the tank ends 5.08~K too cold.}
\summary{Outlet below the front $\to$ smearing drains heat.}
\src{Thesis, Section 5.1.}
\end{qabox}

\begin{qabox}{Why not simply use negative diffusion to sharpen the front?}
\origin{predicted}{your thesis: everything you used can be asked}
\core{Negative diffusion is ill posed: small wiggles grow fastest.}
\explain It does steepen the front, but it also amplifies every disturbance at the grid scale. In my
1D test a disturbance of 1~mK left the range of the data by more than 1~K after 21 steps. So I keep
the useful part and stop the dangerous part with a limiter.
\derive
\begin{align*}
T_i^{n+1} &= T_i^n + \nu\,(T_{i+1}^n - 2T_i^n + T_{i-1}^n) \why{explicit diffusion step} \\
T_i^n &= G^n e^{\mathrm{i}k i\Delta x} \why{one Fourier mode} \\
G &= 1 + \nu\,(2\cos k\Delta x - 2) = 1 - 4\nu\sin^2(k\Delta x/2) \why{insert the mode} \\
G &= 1 + 4\lvert\nu\rvert \why{shortest mode $k\Delta x = \pi$, $\nu<0$}
\end{align*}
$\lvert G\rvert > 1$, so the shortest mode grows in every step.
\summary{Negative diffusion grows wiggles. Keep the steepening, limit the rest.}
\src{Thesis, Section 4.2.2, Eq.~(4.13), Figure 4.1.}
\end{qabox}

\begin{qabox}{Explain your limiter.}
\origin{predicted}{your thesis: everything you used can be asked}
\core{Each face may move heat up the gradient only as far as both cells stay inside their local bounds.}
\explain For every cell I take the smallest and the largest value of the cell and its neighbours.
Across a face, heat moves from the colder to the warmer cell. The amount is the smaller of the two
headrooms, times a strength $s$. This is the limiter idea of flux-corrected transport, by Zalesak.
\derive
\begin{align*}
Q_h^{+} &= C_h\,(T_h^{\max} - T_h),\quad Q_c^{-} = C_c\,(T_c - T_c^{\min}) \why{headrooms, Eq.~(4.15)} \\
\Delta q_f &= s\,\min(Q_h^{+},Q_c^{-}) \why{limited exchange, Eq.~(4.16)} \\
C_i\,\Delta T_i &\le n_f\,s\,Q_i^{+} \why{at most $n_f$ faces bring heat} \\
T_i + \Delta T_i &\le T_i^{\max} \quad\text{if}\quad s \le 1/n_f = 1/6 \why{Eq.~(4.17)}
\end{align*}
\summary{Smaller headroom times $s$. Conservative, no new extremum, safe for $s \le 1/6$.}
\src{Thesis, Section 4.2.3; Appendix D.4.}
\end{qabox}

\begin{qabox}{Why the direction weight?}
\origin{predicted}{your thesis: everything you used can be asked}
\core{The front is horizontal, so the correction should act across it and not along it.}
\explain Each face gets the weight $\cos^2$ of the angle between its link and the vertical. Faces
across the front get one, faces along it get zero. Without the weight the correction roughens the
front sideways. This weight is exactly the face coefficient of a diffusion that acts in one
direction only.
\derive
\begin{align*}
K_f &= \hat{\mathbf{e}}^{\mathsf{T}}\mathbf{K}\,\hat{\mathbf{e}} \why{two-point face coefficient} \\
    &= K\,(\hat{\mathbf{e}}\cdot\hat{\mathbf{n}})^2 \why{rank-one tensor $\mathbf{K} = K\hat{\mathbf{n}}\hat{\mathbf{n}}^{\mathsf{T}}$} \\
    &= K\,w^{\mathrm{dir}}_f \why{Eq.~(4.21)}
\end{align*}
\summary{Act across the front: $\cos^2$ weight = one-direction diffusion.}
\src{Thesis, Section 4.2.4, Eqs.~(4.19), (4.21).}
\end{qabox}

\begin{qabox}{The recorded diffusivity did not help. Why?}
\origin{predicted}{your thesis: everything you used can be asked}
\core{The recorded values are larger than the explicit exchange can apply, so every face hits the cap.}
\explain The library exchanges at most half a cell volume per face and step. The recorded
diffusivity asks for about nine times more. So every face exchanges the maximum, and the field acts
exactly like the constant $\fsd = 0.5$. The replays are bit-identical, which proves it. A shorter
step or an implicit exchange would let the field act.
\derive
\begin{align*}
4 \times \tfrac18 &= \tfrac12 \why{at most four loops of 1/8 of a cell} \\
\mathrm{Di} &= D\,\Delta t/c_s^{2} \le 0.5 \why{the cap as a diffusion number, Eq.~(2.16)}
\end{align*}
\summary{Recorded $D$ above the cap $\to$ acts like a constant.}
\src{Thesis, Section 5.2.1, Figure 5.4.}
\end{qabox}

\begin{qabox}{What happened under the cooled lid?}
\origin{predicted}{your thesis: everything you used can be asked}
\core{On the fluid of the reference, the corrected replay ends 0.06~K from the reference.}
\explain The reference fluid has the specific heat 1006.43~J/(kg\,K). A replay given the specific
heat of water cools four times less for the same wall power, and ends 4.34~K too warm. On the right
fluid the replay is close, and the correction puts the heat back in the right layer.
\nums{Cooling depth 9.62~K. Without correction $-0.66$~K. With correction at $s\approx0.55$:
$+0.06$~K, end profile 0.41~K.}
\derive
\begin{align*}
\Delta T_{\mathrm{cell}} &= -\frac{q\,\Delta t}{\rho\,c_p\,V} \why{wall loss of one cell, Eq.~(4.25)} \\
\frac{\Delta T_{\mathrm{rep}}}{\Delta T_{\mathrm{case}}} &= \frac{c_{\mathrm{case}}}{c_{\mathrm{rep}}} = \frac{1006.43}{4182} = \frac{1}{4.155} \why{same power, two specific heats} \\
h_{\mathrm{eff}} &= G\,h\,\frac{c_{\mathrm{case}}}{c_{\mathrm{rep}}} \why{Eq.~(4.26)}
\end{align*}
\summary{Same fluid on both sides first. Then the correction does the rest.}
\src{Thesis, Sections 6.1--6.2; Appendix D.5.}
\end{qabox}

\begin{qabox}{What could not be repaired, and why?}
\origin{predicted}{your thesis: everything you used can be asked}
\core{While the tank fills, the layers separate too slowly, and no local rule can fix that.}
\explain A local correction only moves heat between neighbours, so it can change the shape of the
front but not carry heat from one layer to another. This is missing transport. A vertical flux read
from the recording supplies that transport. It reduces the error, but the tank then ends too cold.
\nums{Heterogeneity error over the warm inflow 2.25~K $\to$ 0.45~K. End 0.86~K too cold.}
\derive
\begin{align*}
\rho c_p v T &= \rho c_p v\,\bar{T}_{\mathrm{b}} + \rho c_p v\,(T - \bar{T}_{\mathrm{b}}) \why{split $T$ at the band mean} \\
F &= \rho c_p v\,(T - \bar{T}_{\mathrm{b}}) \why{the recorded flux, Eq.~(2.18)} \\
\langle F\rangle_{\mathrm{b}} &= \rho c_p \langle v' T'\rangle_{\mathrm{b}} \why{band mean: the eddy heat flux}
\end{align*}
\summary{Local repair fixes shape. Missing transport needs recorded transport.}
\src{Thesis, Sections 6.5--6.6; Appendix D.2.}
\end{qabox}

\begin{qabox}{How do you know your results are correct?}
\origin{asked}{thesis part, Ulrich (``statistical methods employed, why other methods were not chosen'')}
\core{Every replay is deterministic, and every run changes one setting against a control.}
\explain A rerun gives the same trace, bit for bit. Replay and reference are compared at the same
instants and over the same end window. The energy budget of each case is closed term by term.
\nums{End window 580 to 599~s on both sides.}
\summary{Deterministic, one change per run, same clock, closed budget.}
\src{Thesis, Section 3.9 and Appendix B.5.}
\end{qabox}

\begin{qabox}{Why these error measures? Bias and RMSE?}
\origin{predicted}{your thesis: everything you used can be asked}
\core{The mean temperature says how much heat there is; the profile and the front width say where it is.}
\explain A replay can be right on one and wrong on the other, so I report both. On a trace, the
bias gives the side and the RMSE gives the size. Their identity separates an offset from a wrong
shape.
\derive
\begin{align*}
\mathrm{RMSE}^2 &= \frac1N\sum_k e_k^2 \why{definition, Eq.~(4.3)} \\
 &= \frac1N\sum_k \bigl[(e_k-\mathrm{bias}) + \mathrm{bias}\bigr]^2 \why{add and subtract the bias} \\
 &= \sigma_e^2 + 2\,\mathrm{bias}\cdot 0 + \mathrm{bias}^2 \why{$\sum_k (e_k - \mathrm{bias}) = 0$} \\
 &= \mathrm{bias}^2 + \sigma_e^2
\end{align*}
\summary{How much heat versus where it is. RMSE$^2$ = bias$^2$ + spread$^2$.}
\src{Thesis, Section 4.1.2, Eq.~(4.3).}
\end{qabox}

\begin{qabox}{Why is the flow laminar?}
\origin{predicted}{your thesis: everything you used can be asked}
\core{The jet enters at a Reynolds number of about 99, far below transition.}
\explain So the equations are solved as they stand, with no turbulence model. The case file
confirms that the turbulence model is off. Buoyancy is slightly stronger than inertia, so the jet
bends up within a few inlet widths.
\nums{$Re = UD/\nu \approx 99$ with $U = 0.02$~m/s, $D = 5$~mm. $Ri = 1.84$.}
\summary{$Re \approx 99$: laminar. $Ri = 1.84$: the jet bends up.}
\src{Thesis, Section 3.3, Table 3.2.}
\end{qabox}

\begin{qabox}{Why the Boussinesq approximation?}
\origin{predicted}{your thesis: everything you used can be asked}
\core{The density changes by only 1.5 per cent, so it matters only where gravity acts.}
\explain Density is kept constant except in the buoyancy term. This buoyancy term is exactly the
link that the replay cuts: the replay keeps the recorded motion, so the temperature can no longer
change the velocity.
\derive
\begin{align*}
\rho\,\mathbf{g} &\approx \rho_0\,\mathbf{g}\bigl(1 - \beta(T - T_0)\bigr) \why{linear in $T$ for small changes} \\
&= \rho_0\,\mathbf{g} - \rho_0\,\beta\,(T-T_0)\,\mathbf{g} \why{second term = buoyancy}
\end{align*}
\summary{Constant density, except where gravity sees it. The replay cuts this term.}
\src{Thesis, Section 2.1, Eqs.~(2.1)--(2.3).}
\end{qabox}

\begin{qabox}{What is future work?}
\origin{asked}{thesis part, Küng committee (``What would you do differently?''), Mara (``How could I extend my work?'')}
\core{Let the recorded diffusivity act, split the conservation constraint, and switch databases by state.}
\explain An implicit exchange or a shorter step would let the recorded diffusivity work. A balance
per compartment would stop the constraint from putting heat back at the front. Several recordings,
chosen by the replayed temperature, could cover more than one regime.
\summary{Implicit exchange, compartments, several databases.}
\src{Thesis, Section 9.3.}
\end{qabox}

\begin{qabox}{Which papers influenced you most?}
\origin{asked}{thesis part, Mara (``Which was the paper your thesis was most influenced by?'')}
\core{Lichtenegger and Pirker (2016) and Pirker et al.\ (2018) for rCFD, and Zalesak (1979) for the limiter.}
\explain The first introduced recurrence CFD and the recurrence matrix. The second describes the
face exchange and the conservation constraint. Zalesak gives the local bounds of flux-corrected
transport that my limiter uses.
\summary{rCFD: Lichtenegger, Pirker. Limiter: Zalesak.}
\src{Thesis, Chapter 2 and Section 4.2.}
\end{qabox}
